\documentclass[a4paper,12pt]{article}
\usepackage[T2A]{fontenc}
\usepackage[utf8]{inputenc}
\usepackage[russian]{babel}
\usepackage{amsmath,amssymb}
\usepackage{hyperref}
\usepackage{geometry}
\geometry{top=2cm, bottom=2cm, left=2.5cm, right=2.5cm}

\title{Конспект лекции 4: \\ Рекуррентные нейронные сети}
\author{Липкин С.М.}
\date{}

\begin{document}
\maketitle

\section{RNN: базовые понятия}

Рекуррентные нейронные сети (RNN) — класс сетей для обработки последовательных данных (тексты, речь, временные ряды, видео, ДНК). Ключевые особенности: внутренняя память $h_t$, разделяемые веса, обработка произвольной длины.

\subsection{Математика Vanilla RNN}
Скрытое состояние:
\[
h_t = \tanh(W_{xh} x_t + W_{hh} h_{t-1} + b_h)
\]
Выход:
\[
y_t = W_{hy} h_t + b_y
\]
Параметры $W, b$ разделяются на всех временных шагах.

\section{Обучение: BPTT}

Backpropagation Through Time:
\begin{enumerate}
    \item Forward pass — обработка шаг за шагом.
    \item Loss $L = \sum_{t=1}^{T} L_t$.
    \item Backward — обратное распространение от $T$ к $1$.
    \item Update — обновление весов по сумме градиентов.
\end{enumerate}
Для длинных последовательностей используется Truncated BPTT (window $k \approx 20$--50).

\section{Проблемы градиентов}

\subsection{Исчезающий градиент}
\[
\frac{\partial h_t}{\partial h_j} = \prod_{k=j+1}^{t} W_{hh}^{T} \cdot \operatorname{diag}(f'(z_k))
\]
Если $\max |\lambda(W_{hh})| < 1$, произведение $\to 0$ — веса не обновляются, модель забывает контекст.

\subsection{Взрывающийся градиент}
Если $\max |\lambda(W_{hh})| > 1$, градиент экспоненциально растёт $\to$ NaN, нестабильность.
Решение: \textbf{Gradient Clipping}:
\[
g = \text{threshold} \cdot \frac{g}{\|g\|}, \quad \text{threshold} \in [1, 5]
\]

\section{LSTM (Long Short-Term Memory)}

LSTM решает проблему градиентов через аддитивное обновление cell state $C_t$.

\subsection{Вентили}
\begin{align*}
f_t &= \sigma(W_f [h_{t-1}, x_t] + b_f) \quad \text{(forget gate)} \\
i_t &= \sigma(W_i [h_{t-1}, x_t] + b_i) \quad \text{(input gate)} \\
\tilde{C}_t &= \tanh(W_C [h_{t-1}, x_t] + b_C) \\
o_t &= \sigma(W_o [h_{t-1}, x_t] + b_o) \quad \text{(output gate)}
\end{align*}

\subsection{Обновление состояния}
\[
C_t = f_t \odot C_{t-1} + i_t \odot \tilde{C}_t, \qquad
h_t = o_t \odot \tanh(C_t)
\]
Градиент протекает через $C_t$ — исчезающий градиент не возникает.

\section{GRU (Gated Recurrent Unit)}

Упрощение LSTM: два вентиля (update $z_t$, reset $r_t$), единое состояние $h_t$. Меньше параметров, быстрее обучение, меньше переобучение. На небольших датасетах часто не уступает LSTM.

\section{Архитектуры seq2seq}

\begin{itemize}
    \item \textbf{Many-to-One}: классификация последовательности (sentiment analysis).
    \item \textbf{One-to-Many}: генерация (image captioning).
    \item \textbf{Many-to-Many}: машинный перевод, распознавание речи.
    \item \textbf{Encoder-Decoder}: энкодер сжимает вход в контекстный вектор, декодер генерирует выход.
\end{itemize}

\section{Продвинутые архитектуры}

\subsection{Bidirectional RNN}
Два RNN (forward + backward), конкатенация состояний.

\subsection{Deep RNN}
Несколько слоёв RNN для иерархии признаков.

\subsection{Attention}
Взвешивание важности частей последовательности:
\[
\alpha_i = \text{softmax}(e_i), \quad \text{context} = \sum \alpha_i h_i
\]
Attention $\to$ Transformers.

\section{Ограничения RNN и современные тенденции}

Ограничения:
\begin{enumerate}
    \item Последовательная обработка — нет параллелизма, медленно.
    \item Потеря долгосрочного контекста (даже с LSTM/GRU).
\end{enumerate}
Тенденции:
\begin{itemize}
    \item \textbf{Transformers}: полный параллелизм, self-attention (BERT, GPT).
    \item RNN остаются полезны для stream/edge/real-time задач.
\end{itemize}

\end{document}
